% ch4_d (书页 187-192 / p202-p207)
%--C-TAIL--
% ---- 书页 187 / p202 ----（承接 4.8 节“替代理论”：前块止于书页 186 段末 “…where indices $a,b$ run from 0 to $d+3$.”，前句完整收束，本 tail 直接由式 (4.143) 起，无需 %JOIN%）
\begin{equation*}
ds^{2} = G_{ab}\,dX^{a}dX^{b} = g_{\mu\nu}(x)\,dx^{\mu}dx^{\nu} + b^{2}(x)\gamma_{ij}(y)\,dy^{i}dy^{j},
\tag{4.143}
\end{equation*}
其中 $x^{\mu}$ 是四维时空中的坐标，而 $y^{i}$ 是额外维流形上的坐标；这个流形取为一个具有度规 $\gamma_{ij}$ 的最大对称空间。当然，额外维的几何实际上是某种动力学的东西，本应通过求解完整的运动方程来确定，但我们打算把 (4.143) 当作一个简化拟设（ansatz）。（在更完整的处理中，我们会把紧化几何的动力学模式按傅里叶级数展开，并证明我们眼下忽略的那些模式比我们所选择考察的整体尺寸模式质量更大。）作用量就是 $(4+d)$ 维的 Hilbert 作用量再加上一个物质项：
\begin{equation*}
S = \int d^{4+d}X\sqrt{-G}\left(\frac{1}{16\pi G_{4+d}}R[G_{ab}] + \widehat{\mathcal{L}}_M\right),
\tag{4.144}
\end{equation*}
其中 $\sqrt{-G}$ 是 $G_{ab}$ 的行列式取负号后的平方根，$R[G_{ab}]$ 是 $G_{ab}$ 的 Ricci 标量曲率，而 $\widehat{\mathcal{L}}_M$ 是提出度规行列式因子之后的物质拉格朗日密度。

第一步是对作用量 (4.144) 作维数约化。我们的意思是切实地对额外维完成积分，这之所以可能，是因为我们已经假定额外维的尺度因子 $b$ 与 $y^{i}$ 无关。于是我们可以把一切都用 $g_{\mu\nu}$、$\gamma_{ij}$ 和 $b(x)$ 表出，对额外维积分，从而得到一个有效的四维理论。由度规 (4.143) 可得
\begin{equation*}
\sqrt{-G} = b^{d}\sqrt{-g}\sqrt{\gamma},
\tag{4.145}
\end{equation*}
并且我们可以对这个度规计算标量曲率，得到
\begin{align*}
R[G_{ab}] = R[g_{\mu\nu}] &+ b^{-2}R[\gamma_{ij}] - 2db^{-1}g^{\mu\sigma}\nabla_{\mu}\nabla_{\sigma}b\\
&- d(d-1)b^{-2}g^{\mu\sigma}(\nabla_{\mu}b)(\nabla_{\sigma}b),
\tag{4.146}
\end{align*}
其中 $\nabla_{\mu}$ 是与四维度规 $g_{\mu\nu}$ 相联系的协变导数。我们用 $\mathcal{V}$ 表示 $b=1$ 时额外维的体积；它由下式给出：
\begin{equation*}
\mathcal{V} = \int d^{d}y\sqrt{\gamma}.
\tag{4.147}
\end{equation*}
四维的 Newton 引力常数 $G_{4}$ 由作用量中标量曲率的系数确定；我们发现 $G_{4}$ 与其高维对应物的关系为
\begin{equation*}
\frac{1}{16\pi G_{4}} = \frac{\mathcal{V}}{16\pi G_{4+d}}.
\tag{4.148}
\end{equation*}
于是我们得到
% ---- 书页 188 / p203 ----
\begin{align*}
S = \int d^{4}x\sqrt{-g}\,\bigg\{ \frac{1}{16\pi G_{4}}\Big[ b^{d}R[g_{\mu\nu}] &+ d(d-1)b^{d-2}g^{\mu\nu}(\nabla_{\mu}b)(\nabla_{\nu}b)\\
&+ d(d-1)\kappa b^{d-2}\Big] + \mathcal{V}b^{d}\widehat{\mathcal{L}}_M\,\bigg\},
\tag{4.149}
\end{align*}
其中为了方便我们作了分部积分，并且引入了 $\gamma_{ij}$ 的曲率参数 $\kappa$，由下式给出：
\begin{equation*}
\kappa = \frac{R[\gamma_{ij}]}{d(d-1)}.
\tag{4.150}
\end{equation*}

与式 (4.117)–(4.120) 比较可见，这个维数约化后的作用量恰恰就是一个标量-张量理论的作用量；额外维的尺寸扮演着标量场的角色。因此，我们可以通过一次变量代换加一个共形变换，让它看起来更合乎常规：
\begin{gather*}
\beta(x) = \ln b,\\
\tilde{g}_{\mu\nu} = e^{d\beta}g_{\mu\nu},
\tag{4.151}
\end{gather*}
这便把约化后的作用量变成了 Einstein 框架中一个与引力相耦合的标量场的作用量。按照我们讨论标量-张量理论时所概述的同一流程，就得到
\begin{align*}
S = \int d^{4}x\sqrt{-\tilde{g}}\,\bigg\{ \frac{1}{16\pi G_{4}}\Big[ R[\tilde{g}_{\mu\nu}] &- \frac{1}{2}d(d+2)\tilde{g}^{\mu\nu}(\widetilde{\nabla}_{\mu}\beta)(\widetilde{\nabla}_{\nu}\beta)\\
&+ d(d-1)\kappa e^{(d+2)\beta}\Big] + \mathcal{V}e^{-d\beta}\widehat{\mathcal{L}}_M\,\bigg\},
\tag{4.152}
\end{align*}
其中我们略去了全导数项。

为了把 $\beta$ 变成正则归一化的标量场，我们作最后一次变量代换：
\begin{equation*}
\phi = \sqrt{\frac{d(d+2)}{2}}\,\bar{m}_P\beta,
\tag{4.153}
\end{equation*}
其中约化 Planck 质量为 $\bar{m}_P=(8\pi G_{4})^{-1/2}$。于是我们得到
\begin{align*}
S = \int d^{4}x\sqrt{-\tilde{g}}\,\bigg\{ \frac{1}{16\pi G_{4}}R[\tilde{g}_{\mu\nu}] &- \frac{1}{2}\tilde{g}^{\mu\nu}(\widetilde{\nabla}_{\mu}\phi)(\widetilde{\nabla}_{\nu}\phi)\\
&+ \frac{1}{2}\kappa d(d-1)\bar{m}_P^{2}e^{-\sqrt{2(d+2)}/d\,\phi/\bar{m}_P} + \mathcal{V}e^{-\sqrt{2d/(d+2)}\,\phi/\bar{m}_P}\widehat{\mathcal{L}}_M\,\bigg\}.
\tag{4.154}
\end{align*}

% ---- 书页 189 / p204 ----
标量 $\phi$ 被称为\textbf{伸缩子}（dilaton）或\textbf{径子}（radion），它刻画了额外维流形的尺寸。

式 (4.154) 中的后两项代表（负的）势 $V(\phi)$。如果忽略物质项 $\widehat{\mathcal{L}}_M$，伸缩子的行为将只取决于 $\kappa$ 的符号。如果额外维流形是平直的（$\kappa=0$），势消失，我们就只有一个无质量标量场；这种可能性与我们前面提到的标量-张量理论所受的实验限制相冲突。如果有曲率（$\kappa\neq 0$），势就没有极小值：$\kappa>0$ 时场会滚向 $-\infty$，而 $\kappa<0$ 时场会滚向 $+\infty$。但 $\phi\propto\ln b$，所以这意味着额外维的尺度因子 $b(x)$ 要么收缩到零，要么变得任意大，无论哪种情形都让稳定额外维的希望化为泡影。不过，通过选取一个恰当的物质拉格朗日量，以及额外维中一个恰当的场位形，稳定性是可以实现的。

现在让我们转向另一类替代理论：那些拉格朗日量含有度规导数的高于二阶的项的理论。我们可以想象一个如下形式的作用量
\begin{equation*}
S = \int d^{n}x\sqrt{-g}\left(R + \alpha_{1}R^{2} + \alpha_{2}R_{\mu\nu}R^{\mu\nu} + \alpha_{3}g^{\mu\nu}\nabla_{\mu}R\nabla_{\nu}R + \cdots\right),
\tag{4.155}
\end{equation*}
其中各个 $\alpha$ 是耦合常数，省略号代表我们用曲率张量、它的缩并以及它的导数所能造出的所有其他标量。传统上，这类项一直被忽略，理由相当充分：它们只是把一个已经在美学上令人愉悦、在经验上成功的理论弄得更复杂。除此之外，从经典的角度看还有一个更实质性的反对意见。在常规形式下，Einstein 方程为度规导致一个适定的初值问题：在初始时刻给定的坐标和动量，可以用来预言未来的演化。而若带有高阶导数项，我们不仅需要那些数据，还需要动量的若干阶导数；理论的性质就被急剧改变了。

不过，也有很好的理由去考虑这类附加项。正如我们在对量子引力的简短讨论中提到过的，广义相对论自洽量子化的技术障碍之一是这个理论不可重整化：把高阶量子效应包括进来，就会得到无穷大的答案。而取高阶拉格朗日项的恰当组合，你实际上可以让理论变得可重整化，这为构造一个自洽的量子理论带来了一些希望。\footnote{例如可参见 K.S. Stelle, \emph{Phys. Rev.} \textbf{D16}, 953 (1977).}不幸的是，事实证明可重整化要付出的代价太高：这些模型一般都含有负能量的场激发（幽灵，ghosts）。这样一来，所谓的真空态（空无一物的空间）对于衰变成正、负能量模式而言将是不稳定的，这既与经验事实不一致，也与我们理论上的成见相矛盾。

尽管如此，目前流行的观点是：GR 是一个在 Planck 标度以下的能量处有效的有效理论，我们实际上应当把所有可能的高阶项都包括进来；
% ---- 书页 190 / p205 ----（承接上句：“...include all of the pos-” 续接 “sible higher-order terms...”）
只不过它们会被 Planck 标度的适当幂次压低，正如我们在 4.7 节讨论等效原理时所论证的那样。因此，只有当曲率的特征长度逼近 Planck 标度时，这些项才会变得重要，而这远离任何说得通的实验。所以，高阶项在原则上有趣，在实践中却不然。另一方面，类似的推理会让我们预期一个巨大的真空能项，因为它比 Hilbert 作用量的阶更低，而我们知道这并不是事实；所以我们应该保持开放的心态。

作为广义相对论的最后一个替代方案，我们应当提一提这样一种可能性：联络其实并非从度规导出，而是作为一个基本场独立存在。习题中有一道题要你证明：可以考察广义相对论的常规作用量，但把它视为度规 $g_{\mu\nu}$ 和一个无挠联络 $\Gamma^{\lambda}_{\rho\sigma}$ 二者的函数，而对这样一个作用量关于联络作变分所得的运动方程，意味着 $\Gamma^{\lambda}_{\rho\sigma}$ 实际上就是与 $g_{\mu\nu}$ 相联系的 Christoffel 联络。我们也可以放弃联络必须无挠的要求，这时挠率张量可能带来额外的传播自由度。这类理论之所以没有受到多少关注，基本原因很简单：挠率本身是一个张量，没有任何东西把它与其他非引力的张量场区分开来。因此，考虑无挠联络的理论（它们给出 GR）再加上任意多个我们爱怎么命名就怎么命名的张量场，其实并没有损失任何一般性。当我们考虑放弃度规相容性要求时，类似的考虑同样适用——任何联络都可以写成一个度规相容的联络加上一个张量修正，于是任何这类理论都等价于 GR 加上额外的张量场，而后者实在不配称作“广义相对论的替代理论”。

\section{习题}

% 习题编号照原书
\textbf{1.}\quad 弯曲空间中电磁学的拉格朗日密度为
\begin{equation*}
\mathcal{L} = \sqrt{-g}\left(-\frac{1}{4}F^{\mu\nu}F_{\mu\nu} + A_{\mu}J^{\mu}\right),
\tag{4.156}
\end{equation*}
其中 $J^{\mu}$ 是守恒流。

\textbf{(a)}\quad 通过对度规作泛函微分，导出能动张量。你可以假定 $A_{\mu}J^{\mu}$ 项对能动张量没有贡献。

\textbf{(b)}\quad 考虑在拉格朗日量中加进一个新项，
\begin{equation*}
\mathcal{L}' = \beta R^{\mu\nu}g^{\rho\sigma}F_{\mu\rho}F_{\nu\sigma}.
\end{equation*}
在这一项存在时，麦克斯韦方程组如何改变？Einstein 方程呢？流还守恒吗？

\textbf{2.}\quad 我们已经展示过如何通过对 Hilbert 作用量关于度规作变分来导出 Einstein 方程。它们也可以这样导出：把度规和联络当作独立的自由度，分别对它们作变分；这就是所谓的
% ---- 书页 191 / p206 ----（承接上句：“...this is known” 续接 “as the Palatini formalism.”）
\textbf{Palatini 形式}（Palatini formalism）。也就是说，我们考虑作用量
\begin{equation*}
S = \int d^{4}x\sqrt{-g}\,g^{\mu\nu}R_{\mu\nu}(\Gamma),
\end{equation*}
其中 Ricci 张量被视为纯粹由联络构造而成，不使用度规。对度规的变分给出通常的 Einstein 方程，只不过其中的 Ricci 张量是由一个与度规没有任何先验关系的联络构造出来的。一开始就设想联络是对称的（无挠的），试证明：对这个作用量关于联络系数作变分，将导致联络必须度规相容的要求，也就是说，联络是 Christoffel 联络。记住，把矢量的协变散度的积分与矢量在边界上的积分联系起来的 Stokes 定理，对一般的协变导数并不成立。最好的策略是把联络系数写成 Christoffel 符号 $\widetilde{\Gamma}^{\lambda}_{\mu\nu}$ 与一个张量 $C^{\lambda}_{\mu\nu}$ 之和，
\begin{equation*}
\Gamma^{\lambda}_{\mu\nu} = \widetilde{\Gamma}^{\lambda}_{\mu\nu} + C^{\lambda}_{\mu\nu},
\end{equation*}
然后证明 $C^{\lambda}_{\mu\nu}$ 必定为零。

\textbf{3.}\quad 流形 $M$ 上的四维 $\delta$ 函数定义为
\begin{equation*}
\int_{M}F(x^{\mu})\left[\frac{\delta^{(4)}(x^{\sigma}-y^{\sigma})}{\sqrt{-g}}\right]\sqrt{-g}\,d^{4}x = F(y^{\sigma}),
\tag{4.157}
\end{equation*}
其中 $F(x^{\mu})$ 为任意函数。另一方面，无压强理想流体（尘埃）的能动张量为
\begin{equation*}
T^{\mu\nu} = \rho U^{\mu}U^{\nu},
\tag{4.158}
\end{equation*}
其中 $\rho$ 是能量密度，$U^{\mu}$ 是四速度。考虑这样一个流体：它由单独一个沿世界线 $x^{\mu}(\tau)$ 运动的粒子构成，$\tau$ 为固有时。这个流体的能动张量便由下式给出：
\begin{equation*}
T^{\mu\nu}(y^{\sigma}) = m\int_{M}\left[\frac{\delta^{(4)}(y^{\sigma}-x^{\sigma}(\tau))}{\sqrt{-g}}\right]\frac{dx^{\mu}}{d\tau}\frac{dx^{\nu}}{d\tau}\,d\tau,
\tag{4.159}
\end{equation*}
其中 $m$ 是粒子的静质量。试证明：能动张量的协变守恒，$\nabla_{\mu}T^{\mu\nu}=0$，蕴含 $x^{\mu}(\tau)$ 满足测地线方程。

\textbf{4.}\quad 试证明电磁学与标量场论的能动张量满足主能量条件，从而也满足弱、类光以及类光主条件。试证明它们还满足 $w\geq -1$。

\textbf{5.}\quad 如果存在一个与类空超曲面正交的类时 Killing 矢量，我们就说一个时空是静态的。（更多的讨论，包括 Raychaudhuri 方程的定义，见附录 D 与附录 F。）

\textbf{(a)}\quad 一般地说，如果矢量场 $v^{\mu}$ 与由 $f=$ 常数定义的一族超曲面正交，那么我们可以把这个矢量写成 $v_{\mu}=h\nabla_{\mu}f$（这里 $f$ 和 $h$ 都是函数）。试证明这蕴含
\begin{equation*}
v_{[\sigma}\nabla_{\mu}v_{\nu]} = 0.
\end{equation*}

% ---- 书页 192 / p207 ----
\textbf{(b)}\quad 想象我们有一个零压强的理想流体（尘埃），它生成 Einstein 方程的一个解。试证明：只有当流体四速度平行于类时（且超曲面正交）的 Killing 矢量时，度规才可能是静态的。

\textbf{(c)}\quad 利用 Raychaudhuri 方程证明：如果压强为零而能量密度大于零，则 Einstein 方程不存在静态解。

\textbf{6.}\quad 设 $K$ 是一个 Killing 矢量场。试证明：如果度规是 Einstein 方程的真空解，那么势为 $A_{\mu}=K_{\mu}$ 的电磁场求解麦克斯韦方程组。这有点作弊，因为如果电磁场强非零，你就不再处于真空之中；但我们假定场强足够小，不至于显著地影响几何。

% 第4章正文至此结束（原书书页 192 末；书页 193 起为第 5 章）。
